% Section 3.11, from lectures/L03/appendix/b_exercises.md.

\section{Exercises}
\label{c3:sec:exercises}

\begin{aside}
\textbf{How to work these.} Each exercise is tagged with the kinds of work it asks for:
\kinds{Paper} by hand, \kinds{CircuitVerse} in the browser, \kinds{VHDL} in a module of your own.

\textbf{How to check your work.} Every exercise that asks you to write a VHDL module ships a
self-checking testbench under \file{lectures/L03/exercises} in the companion repository. Write your
module in its \file{exercises/<module>} directory, using the entity name and the \textbf{port order}
the exercise specifies, then run it with GHDL; \secref{c2:app:c} has the three commands.

No FPGA board is needed for any exercise. Your job is to get the VHDL right, and the testbench is
how you confirm it. Short answers to the parts you work on paper are in
Appendix~\ref{app:answers}; reference solutions to the VHDL are published in the companion
repository, under \file{lectures/L03/solutions}. Try each exercise yourself before looking at
either.
\end{aside}

% ------------------------------------------------------------------------------------------
\exerciseset{D Latch}
\label{c3:set:1}

\exercise{A D latch}{Paper, CircuitVerse}
\label{c3:ex:1}
Build the first circuit in this book that can \emph{remember} something.

Everything you have built so far has been combinational: the outputs follow the inputs, and the
moment the inputs change the outputs change with them. A latch is the smallest departure from that.
It has a data input \code{D} and a control input \code{enable}, and it has two outputs, \code{Q} and
its complement \code{Qn}. While \code{enable} is high the latch is \textbf{transparent}: \code{Q}
simply follows \code{D}, and it looks like a wire. The interesting part is what happens when
\code{enable} goes low. The latch \textbf{locks}, holding whatever \code{D} was at that instant, and
from then on \code{D} can do what it likes without \code{Q} noticing. That held value is one bit of
memory, and it is built out of nothing but the gates you already know, wired so that the output
feeds back into the input.

\code{Qn} is not decoration. The feedback loop is what makes the circuit remember, and the two
outputs are what feed each other: \code{Q} is computed from \code{Qn} and \code{Qn} from \code{Q}.
Build it and you can see the loop; that is the point of drawing it by hand before you ever meet the
VHDL.

Use the equations from \secref{c3:a2}:
\[ Q = (D' \cdot \bool{enable} + \bool{Qn})' \]
\[ \bool{Qn} = (D \cdot \bool{enable} + Q)' \]

\begin{parts}
  \item Realize the corresponding gate network:
    \begin{itemize}
      \item Draw the network by hand.
      \item Recreate and simulate it in CircuitVerse (\url{https://circuitverse.org/simulator}).
    \end{itemize}
  \item Test the transparent and locked states:
    \begin{itemize}
      \item Set \code{enable = 1}:
        \begin{itemize}
          \item Change \code{D}.
          \item Observe what happens to \code{Q} and \code{Qn}.
        \end{itemize}
      \item Set \code{enable = 0}:
        \begin{itemize}
          \item Change \code{D} again.
          \item Observe what happens to \code{Q} and \code{Qn}.
        \end{itemize}
      \item Explain the difference between the two cases.
    \end{itemize}
  \item Cycle \code{D} and \code{enable} through the following combinations twice:
    \begin{itemize}
      \item \code{00}
      \item \code{01}
      \item \code{10}
      \item \code{11}
    \end{itemize}
    Confirm that you observe both cases:
    \begin{itemize}
      \item The latch locks after being transparent with \code{D = 1}.
      \item The latch locks after being transparent with \code{D = 0}.
    \end{itemize}
    Explain how the value of \code{Q} depends on the value of \code{D} at the instant \code{enable}
    changes from \code{1} to \code{0}.
\end{parts}

% ------------------------------------------------------------------------------------------
\exerciseset{D Flip-Flop}
\label{c3:set:2}

\exercise{A D flip-flop from two latches}{Paper, CircuitVerse}
\label{c3:ex:2}
Extend the D latch from Exercise~\ref{c3:ex:1} into a D flip-flop, following \secref{c3:a3}.

The flip-flop must have:
\begin{itemize}
  \item Inputs:
    \begin{itemize}
      \item \code{clock}
      \item \code{reset_n}
      \item \code{D}
      \item \code{enable}
    \end{itemize}
  \item Outputs:
    \begin{itemize}
      \item \code{Q}
      \item \code{Qn}
    \end{itemize}
\end{itemize}
The reset signal is:
\begin{itemize}
  \item Asynchronous:
    \begin{itemize}
      \item It takes effect without waiting for a clock edge.
    \end{itemize}
  \item Active-low:
    \begin{itemize}
      \item Setting \code{reset_n = 0} must force \code{Q} to \code{0} immediately.
    \end{itemize}
\end{itemize}
\note[Tip]
\begin{itemize}
  \item The reset has to do two things in each of the two internal latches, not one:
    \begin{itemize}
      \item Feed an active-high \code{reset} (\code{reset_n} through a NOT gate) into the NOR gate
        that produces \code{Q}. This forces \code{Q} low without waiting for a clock edge.
      \item \code{AND} \code{reset_n} into the AND gate that produces that latch's \code{D} and
        \code{enable} product. This frees \code{Qn} to go high.
      \item Gating only the \code{Q} feedback path is not enough: with \code{enable = 1} and
        \code{D = 1} it leaves \code{Qn} low, and a low \code{Qn} is exactly what holds \code{Q}
        high, so the reset would do nothing.
    \end{itemize}
  \item Drive the enable inputs of the two internal latches using:
    \begin{itemize}
      \item \code{clock}
      \item The inverse of \code{clock}
    \end{itemize}
  \item Follow \secref{c3:a3} for how the external \code{enable} input controls whether the
    flip-flop accepts a new value.
\end{itemize}

\begin{parts}
  \item Realize the corresponding gate network:
    \begin{itemize}
      \item Draw the network by hand.
      \item Recreate and simulate it in CircuitVerse.
      \item Set the clock period to \code{1000 ms}.
        \begin{itemize}
          \item This is slow enough to observe the behaviour by eye.
        \end{itemize}
    \end{itemize}
  \item Test the clocked behaviour:
    \begin{itemize}
      \item Change \code{D} during the high phase of the clock.
      \item Change \code{D} during the low phase of the clock.
      \item Observe \code{Q} and \code{Qn}.
      \item Confirm that the outputs change only on the active clock edge, never between clock
        edges.
    \end{itemize}
    Explain how the two internal latches work together to produce edge-triggered behaviour.
  \item Test the asynchronous reset:
    \begin{itemize}
      \item First set \code{Q = 1}.
      \item Assert the reset by setting \code{reset_n = 0}.
      \item Observe whether \code{Q} resets:
        \begin{itemize}
          \item Immediately.
          \item At the next clock edge.
        \end{itemize}
    \end{itemize}
    Explain the result by referring to the difference between an asynchronous signal and a
    synchronous signal.
\end{parts}

% ------------------------------------------------------------------------------------------
\exerciseset{Registers}
\label{c3:set:3}

\exercise{A four-bit register}{VHDL, Paper}
\label{c3:ex:3}
Write a complete VHDL entity named \code{register4}: four flip-flops side by side, sharing one
clock, one reset, and one enable.

This is the exercise where the flip-flop stops being a curiosity and becomes a building block. A
single D flip-flop stores one bit, which is not much use on its own; what a design actually wants is
somewhere to keep a \emph{number}. Put four flip-flops in a row, wire the same clock to all of them,
and you have a four-bit register. Nothing about the individual flip-flop changes. Only the width
does.

The part worth paying attention to is how little the VHDL has to change to say that. You are
writing the same synchronous process template from \secref{c3:a6} that you already used for one bit,
with \code{std_logic_vector(3 downto 0)} in place of \code{std_logic}. There is no loop, no
repetition, no four copies of anything. That is not a shortcut the language is offering you; it is a
fact about the hardware, and part \textbf{c)} below asks you to say why.

The entity must have:

\begin{booktable}{@{}p{4em}p{4.8em}p{14.6em}L@{}}
  \thead{Port} & \thead{Direction} & \thead{Type} & \thead{Description} \\
  \midrule
  \code{clock} & in & \code{std_logic} & System clock. \\
  \code{reset_n} & in & \code{std_logic} & Asynchronous, active-low. Clears the register to
    \code{"0000"}. \\
  \code{enable} & in & \code{std_logic} & \code{'0'} retains the current value; \code{'1'} captures
    \code{d} on the next rising edge. \\
  \code{d} & in & \code{std_logic_vector(3 downto 0)} & The value to capture. \\
  \code{q} & out & \code{std_logic_vector(3 downto 0)} & The stored value. \\
\end{booktable}

\noindent Write the process out yourself rather than copying the worked D flip-flop example,
\file{d_flip_flop.vhd} (\secref{c3:a8}). Typing it is how the template stops being something you
recognize and starts being something you can produce from memory, and you will be writing it in
every remaining chapter of this book.

\begin{parts}
  \item Write the entity and the synchronous process, adapting the template from a single bit to a
    four-bit vector.
  \item Check it with the testbench, as described at the start of these exercises.
  \item Explain why the same template scales from one flip-flop to four without any structural
    change, while a four-bit \emph{counter} could not simply be four one-bit counters side by side.
    What is different about the two cases?
\end{parts}

\modulefigure{register4}

\selfcheck name your entity \code{register4}, with ports \code{clock}, \code{reset_n},
\code{enable}, \code{d} (\code{std_logic_vector(3 downto 0)}) and \code{q}
(\code{std_logic_vector(3 downto 0)}), declared in that order; its testbench is in
\file{exercises/register4}.

% ------------------------------------------------------------------------------------------
\exerciseset{Edge Detection}
\label{c3:set:4}

\exercise{An edge-detected toggle}{Paper, CircuitVerse}
\label{c3:ex:4}
Construct a circuit that toggles an LED once for every sampled rising edge of a button signal.

Follow:
\begin{itemize}
  \item \secref{c3:a7}
  \item \secref{c3:a9}
\end{itemize}
The circuit must have:
\begin{itemize}
  \item Inputs:
    \begin{itemize}
      \item \code{button}
      \item \code{clock}
      \item \code{reset}
    \end{itemize}
  \item Output:
    \begin{itemize}
      \item \code{led}
    \end{itemize}
\end{itemize}
For this exercise, \code{reset} is active-high, unlike the reset used in the worked example. It is
still \textbf{asynchronous}, like every reset in this book: assert it and \code{led} clears
immediately, without waiting for a clock edge. In the template from \secref{c3:a6} that means the
reset branch stays outside the \code{rising_edge(clock)} test, and only the polarity of the test
changes.

\note[Tip] Use two flip-flops:
\begin{itemize}
  \item One flip-flop stores the previous sampled value of \code{button}.
    \begin{itemize}
      \item Use this value to detect a rising edge.
    \end{itemize}
  \item One flip-flop stores the current state of \code{led}.
    \begin{itemize}
      \item Toggle it only when a rising edge is detected.
    \end{itemize}
\end{itemize}

\begin{parts}
  \item Realize the corresponding gate network:
    \begin{itemize}
      \item Draw the network by hand.
      \item Recreate and simulate it in CircuitVerse.
      \item Set the clock period to \code{1000 ms}.
    \end{itemize}
  \item Test the edge detector:
    \begin{itemize}
      \item Hold \code{button} high for several consecutive clock cycles.
      \item Confirm that \code{led} toggles exactly once.
      \item Confirm that it does not toggle once per clock cycle.
    \end{itemize}
    Explain which stored signal prevents the circuit from repeatedly detecting the same high level
    as a new rising edge.
\end{parts}

\exercise{The toggle in VHDL}{VHDL}
\label{c3:ex:5}
Implement the circuit from Exercise~\ref{c3:ex:4} in VHDL.

\begin{parts}
  \item Create an entity named \code{led_toggle_single} with:

\begin{booktable}{@{}p{4em}p{4.8em}p{5.2em}L@{}}
  \thead{Port} & \thead{Direction} & \thead{Type} & \thead{Description} \\
  \midrule
  \code{clock} & in & \code{std_logic} & System clock. \\
  \code{reset} & in & \code{std_logic} & Asynchronous, and \textbf{active-high}, unlike every other
    module in this book. Assert it and \code{led} clears without waiting for a clock edge, so the
    reset branch belongs outside the \code{rising_edge(clock)} test. \\
  \code{button} & in & \code{std_logic} & Toggles \code{led} once per press. \\
  \code{led} & out & \code{std_logic} & The LED state. \\
\end{booktable}

    Note the single-bit types: the worked example uses \code{std_logic_vector} ports for two buttons
    and two LEDs, and this one drives a single pair.

    Do not copy \file{led_toggle.vhd} (\secref{c3:a9}) directly. Instead:
    \begin{itemize}
      \item Re-derive the two clocked processes yourself:
        \begin{itemize}
          \item One process stores the previous value of \code{button}.
          \item One process stores and toggles the current value of \code{led}.
        \end{itemize}
      \item Derive the edge-detection equation for:
        \begin{itemize}
          \item A rising edge.
          \item An active-high button.
        \end{itemize}
      \item Use the rising-edge equation from \secref{c3:a7}:
        \[ \bool{edge} = \bool{current} \cdot \bool{previous}' \]
    \end{itemize}
    Do not use the falling-edge equation from the worked example, which uses an active-low button.
  \item Verify the design with its testbench (see the aside at the start of these exercises).

    The testbench drives the same sequence you would exercise by hand on a board:
    \begin{itemize}
      \item It asserts \code{reset} and checks \code{led} is cleared.
      \item It raises \code{button} and checks \code{led} toggles exactly once.
      \item It \textbf{holds \code{button} high} for several more clock cycles and checks, after
        every one of them, that \code{led} does \emph{not} toggle again.
      \item It releases and presses again, and checks the second press toggles \code{led} once more:
        back off, since the first press had turned it on.
      \item It releases and presses a third time, so the LED is lit when the reset arrives.
      \item It asserts \code{reset} with \code{led} lit and \textbf{no clock edge following}, and
        checks \code{led} clears anyway. This is the one that fails if you put the reset inside the
        \code{rising_edge(clock)} branch.
    \end{itemize}
    That third check is the one that fails if your edge detector is wrong:
    \begin{itemize}
      \item An \code{led} that toggles on every clock cycle while the button is held means you drove
        the toggle from \code{button}'s \textbf{level} rather than from its edge.
      \item Re-read the equation above: \code{edge} is high only when \code{button} is high
        \emph{now} and its stored previous value is low.
    \end{itemize}
\end{parts}

\modulefigure{led_toggle_single}

\selfcheck name your entity \code{led_toggle_single}, with ports \code{clock}, \code{reset},
\code{button} (in) and \code{led} (out), declared in that order; its testbench is in
\file{exercises/led_toggle_single}. Note \code{reset} is \textbf{active-high} here, unlike the
worked example.
